\documentclass[11pt]{amsart}

\usepackage[T1]{fontenc}
\usepackage{lmodern}
\usepackage{microtype}
\usepackage{mathtools,amssymb,amsthm}
\usepackage[hidelinks]{hyperref}

\hypersetup{
  pdftitle={The 3-cut complex of the circular ladder is a wedge of spheres in degree 2n-4, for n at least 5}
}

\newtheorem{theorem}{Theorem}
\newtheorem{proposition}[theorem]{Proposition}
\newtheorem{lemma}[theorem]{Lemma}
\theoremstyle{definition}
\newtheorem{conjecture}[theorem]{Conjecture}

\DeclareMathOperator{\rank}{rank}
\newcommand{\CL}{\mathrm{CL}}
\newcommand{\bx}{\mathbin{\square}}
\newcommand{\Dl}{\Delta}
\newcommand{\Zm}{\mathbb{Z}}

\title[The 3-cut complex of the circular ladder]
{The 3-cut complex of the circular ladder is a wedge of spheres
in degree $2n-4$, for $n\ge5$}

\date{}

\subjclass[2020]{05E45, 55U10, 05C76}
\keywords{cut complex, circular ladder, prism graph, Alexander duality, simplicial homology,
wedge of spheres}

\begin{document}

\begin{abstract}
Let $\Dl_k(G)$ be the $k$-cut complex of a graph $G$: the simplicial complex on $V(G)$ whose facets
are the sets $\sigma$ with $|\sigma|=|V(G)|-k$ for which $G[V(G)\setminus\sigma]$ is disconnected.
Chauhan, Shukla and Vinayak conjectured that for $m=2$ and $n\ge 5$ the reduced integral homology of
$\Dl_3(K_m\bx C_n)$ is $\Zm^{2n^2-8n+1}$ in degree $mn-4$ and zero elsewhere. We prove this, and
add that for every $n\ge5$ the complex $\Dl_3(C_n\bx K_2)$ is homotopy equivalent to a wedge of
$2n^2-8n+1$ copies of $S^{2n-4}$. The proof is an integer linear algebra computation on the
Alexander dual, which a parity argument confines to dimension~$3$ for every~$n$, together with three
explicit unimodular minors. The integer $2n^2-8n+1$ and the vanishing of $\widetilde H_i$ for
$i\le 2n-7$ are already in print for this exact graph, as a reduced Euler characteristic and a
codimension-$2$ skeleton, in work of Bayer, Denker, Jeli\'c Milutinovi\'c, Rowlands, Sundaram and
Xue. What we do not find in the literature cited here is the vanishing of the two groups those
results leave free, the absence of torsion, and the homotopy type.
\end{abstract}

\maketitle

\section{The conjecture}\label{sec:conj}

For a graph $G$ on the vertex set $V$ with $|V|=N$ and an integer $k\ge1$, the \emph{$k$-cut complex}
$\Dl_k(G)$ is the simplicial complex on $V$ whose facets are exactly the sets $\sigma\subseteq V$ with
$|\sigma|=N-k$ such that the induced subgraph $G[V\setminus\sigma]$ is \emph{disconnected}
\cite{BayerI}. Equivalently,
\begin{equation}\label{eq:facechar}
 \sigma\in\Dl_k(G)
 \iff
 \text{$V\setminus\sigma$ contains a $k$-set $T$ with $G[T]$ disconnected.}
\end{equation}
This is the non-total complex; the \emph{total} $k$-cut complex, which demands that
$G[V\setminus\sigma]$ have no edge at all, is a different object and is not considered here.

Chauhan, Shukla and Vinayak \cite{CSV} state the following as the last numbered conjecture of their
paper. The e-print carries no printed statement numbers, so we quote the statement in full rather
than by number.

\begin{conjecture}[Chauhan--Shukla--Vinayak]\label{conj:csv}
\begin{enumerate}
\item[(i)] For $m\ge2$ and $n=4$,
 $\ \widetilde H_i\bigl(\Dl_3(K_m\bx C_n)\bigr)=
   \begin{cases}
    \Zm^{2m^2-5m+3} & \text{if } i=mn-5,\\
    \Zm^{2m(m-1)} & \text{if } i=mn-4,\\
    0 & \text{otherwise.}
   \end{cases}$
\item[(ii)] For $m=2$ and $n\ge5$,
 $\ \widetilde H_i\bigl(\Dl_3(K_m\bx C_n)\bigr)=
   \begin{cases}
    \Zm^{2n^2-8n+1} & \text{if } i=mn-4,\\
    0 & \text{otherwise.}
   \end{cases}$
\end{enumerate}
\end{conjecture}

Part~(ii) is the statement settled here. With $m=2$ it reads: for every $n\ge5$,
\begin{equation}\label{eq:target}
 \widetilde H_{2n-4}\bigl(\Dl_3(C_n\bx K_2);\Zm\bigr)\cong\Zm^{2n^2-8n+1},
 \qquad
 \widetilde H_i\bigl(\Dl_3(C_n\bx K_2);\Zm\bigr)=0\ \ (i\ne 2n-4).
\end{equation}
Homology is reduced and integral throughout, as in the source: $\Dl_3$ is connected, so unreduced
$H_0\ne0$ would contradict the clause ``$0$ otherwise''.

\section{The result}\label{sec:result}

The Cartesian product is commutative up to isomorphism, so $K_2\bx C_n\cong C_n\bx K_2$ and
part~(ii) of Conjecture~\ref{conj:csv} is a statement about the latter.
Write $\CL_n:=C_n\bx K_2$, the \emph{circular ladder} or $n$-gonal prism: the vertex set is
$V=\{(r,j): r\in\{0,1\},\ j\in\Zm_n\}$, with edges the \emph{rails} $\{(r,j),(r,j+1)\}$ and the
\emph{rungs} $\{(0,j),(1,j)\}$. So $|V|=2n$, $|E|=3n$, and $\CL_n$ is $3$-regular and triangle-free
of girth~$4$. Write $Q_j:=\{(0,j),(0,j{+}1),(1,j),(1,j{+}1)\}$ for the \emph{rung square} at
column~$j$.

\begin{theorem}\label{thm:main}
Let $n\ge5$. Then
\[
 \widetilde H_i\bigl(\Dl_3(\CL_n);\Zm\bigr)\cong
 \begin{cases}
  \Zm^{2n^2-8n+1} & i=2n-4,\\
  0 & \text{otherwise,}
 \end{cases}
\]
so Conjecture~\ref{conj:csv}(ii) holds. Moreover $\Dl_3(\CL_n)$ is homotopy equivalent to a wedge of
$2n^2-8n+1$ copies of $S^{2n-4}$.
\end{theorem}

Part of the content of Theorem~\ref{thm:main} is published, and is due to Bayer, Denker, Jeli\'c
Milutinovi\'c, Rowlands, Sundaram and Xue. They prove \cite{BayerII} that if $G$ has $N$ vertices,
$k\ge3$, $G$ has no cycle of length $\le k$ and $G$ has a cycle of length $k+1$, then $\Dl_k(G)$
contains a complete codimension-$2$ skeleton and
\begin{equation}\label{eq:euler}
 \widetilde\chi\bigl(\Dl_k(G)\bigr)
 =(-1)^{N-k-1}\Bigl(\tbinom{N-1}{k-1}+|Y_{k+1}(G)|-|Z_k(G)|\Bigr),
\end{equation}
where $Y_{k+1}(G)$ is the set of complements of $(k+1)$-cycles and $Z_k(G)$ the set of complements of
connected $k$-sets. The hypotheses hold verbatim for $G=\CL_n$ and $k=3$: girth $4>3$, and $\CL_n$
has $4$-cycles. Substituting $N=2n$, $|Y_4|=n$ and $|Z_3|=6n$ (Lemma~\ref{lem:fvec}) into
\eqref{eq:euler} gives $\widetilde\chi=\binom{2n-1}{2}+n-6n=2n^2-8n+1$. So the integer
$2n^2-8n+1$ is not new, and neither is the vanishing $\widetilde H_i(\Dl_3(\CL_n))=0$ for
$i\le2n-7$, which is immediate from the complete codimension-$2$ skeleton.

Since $\Dl_3(\CL_n)$ is pure of dimension $2n-4$, those two published facts leave three groups
possibly nonzero --- degrees $2n-6$, $2n-5$, $2n-4$ --- constrained by one linear relation. What
Theorem~\ref{thm:main} adds is therefore the vanishing of $\widetilde H_{2n-5}$ and
$\widetilde H_{2n-6}$, torsion-freeness, which \eqref{eq:euler} cannot see, and the homotopy type.
We have not found those three statements in the literature cited here, and we make no priority
claim: no comprehensive literature search was carried out.

\section{The proof}\label{sec:proof}

\subsection{The Alexander dual is three-dimensional, for every \texorpdfstring{$n$}{n}}

For a simplicial complex $K$ on $V$, $|V|=N$, the Alexander dual is
$K^{*}=\{\sigma\subseteq V: V\setminus\sigma\notin K\}$, and combinatorial Alexander duality
\cite{BjornerTancer} gives
\begin{equation}\label{eq:AD}
 \widetilde H_i(K^{*};\Zm)\cong\widetilde H^{\,N-i-3}(K;\Zm)\qquad(i\in\Zm).
\end{equation}
By \eqref{eq:facechar}, $V\setminus\sigma\notin\Dl_3(G)$ if and only if every $3$-subset of $\sigma$
induces a connected subgraph. Hence
\begin{equation}\label{eq:D}
 D:=\Dl_3(G)^{*}
 =\bigl\{S\subseteq V:\ \text{every $3$-subset of $S$ induces a connected subgraph}\bigr\},
\end{equation}
which is manifestly a simplicial complex: a $3$-subset of a subset of $S$ is a $3$-subset of $S$.

\begin{lemma}\label{lem:fvec}
Let $G$ be a cubic triangle-free graph on $N=2n$ vertices with exactly $q$ induced $4$-cycles. Then
$D$ has $f$-vector
\[
 (f_{-1},f_0,f_1,f_2,f_3,f_4,\dots)=\bigl(1,\ 2n,\ \tbinom{2n}{2},\ 6n,\ q,\ 0,\ 0,\dots\bigr),
\]
so $\dim D=3$ whenever $q>0$. Moreover $q=n$ for $G=\CL_n$ with $n\ge5$, while $q=n+2=6$ for
$G=\CL_4$.
\end{lemma}

\begin{proof}
Since $G$ is triangle-free, a $3$-set induces a connected subgraph precisely when it spans exactly
two edges, that is, when it is an induced path $P_3$: zero or one edge leaves an isolated vertex, and
three edges would be a triangle.

A $2$-set has no $3$-subset, so $f_1=\binom{2n}{2}$. A $3$-set lies in $D$ iff it is an induced
$P_3$, and the number of those is $\sum_{v}\binom{\deg v}{2}=2n\binom32=6n$, with no triangle
correction.

Let $|S|=s\ge4$ with $S\in D$ and put $e=|E(G[S])|$. Each of the $\binom s3$ triples of $S$ spans
exactly $2$ edges, and each edge of $G[S]$ lies in exactly $\binom{s-2}{1}=s-2$ of those triples.
Counting edge-triple incidences twice,
\begin{equation}\label{eq:parity}
 (s-2)\,e=2\binom s3 .
\end{equation}
For $s=5$, \eqref{eq:parity} reads $3e=20$, which has no integer solution: $f_4=0$, and hence
$f_j=0$ for all $j\ge4$ because $D$ is a complex. For $s=4$, \eqref{eq:parity} gives $e=4$; a
triangle-free graph on four vertices with four edges is $C_4$, and $G[S]$ is an induced subgraph, so
$S$ spans an induced $4$-cycle. Conversely every triple of an induced $4$-cycle is an induced $P_3$.
So $f_3=q$.

Finally, the $4$-cycles of $\CL_n$. A cycle of $\CL_n$ uses an even number of rungs. With no rung it
lies in one rail, i.e.\ in a copy of $C_n$, which has a $4$-cycle only if $n=4$; the two rails are
then induced $4$-cycles, and this is the source of the two extra squares at $n=4$. With two rungs, at
columns $j$ and $j'$, the cycle traverses a path of length $a\ge1$ in one rail and a path of length
$b\ge1$ in the other, and $2+a+b=4$ forces $a=b=1$, so $j'=j\pm1$ and the cycle is a rung square
$Q_j$; all $n$ of these are induced. With four or more rungs the cycle has at least eight vertices.
Hence $q=n$ for $n\ge5$ and $q=n+2$ for $n=4$.
\end{proof}

So for every $n\ge5$ the simplicial chain complex of $D$ is
\begin{equation}\label{eq:chain}
 0\longrightarrow \Zm^{n}\xrightarrow{\ \partial_3\ }\Zm^{6n}
 \xrightarrow{\ \partial_2\ }\Zm^{n(2n-1)}\xrightarrow{\ \partial_1\ }\Zm^{2n}
 \xrightarrow{\ \partial_0\ }\Zm\longrightarrow 0,
\end{equation}
a complex of \emph{fixed length}, whose entries are the only thing that grows with $n$; this is why a
statement about all $n$ is here a finite piece of linear algebra rather than an induction.

\subsection{Three unimodular minors}

We use the following standard fact: if an integer matrix $M$ has rank $r$ and some $r\times r$ minor
of $M$ equals $\pm1$, then every invariant factor of $M$ equals $1$, because $d_1d_2\cdots d_r$ is
the gcd of the $r\times r$ minors and $d_1\mid d_2\mid\cdots$. We exhibit such a minor for each of
$\partial_1,\partial_2,\partial_3$.

Fix $n\ge5$. Name the $6n$ triples of $D$ by their centre: for $v=(r,j)$ put
\[
 T_1(v)=\{(r,j{-}1),v,(r,j{+}1)\},\qquad
 T_2(v)=\{(r,j{-}1),v,(1{-}r,j)\},
\]
\[
 T_3(v)=\{(r,j{+}1),v,(1{-}r,j)\}.
\]
Since $G$ is triangle-free, the centre of an induced $P_3$ is determined by its vertex set, so these
$3\cdot 2n=6n$ triples are distinct and exhaust $f_2$.

\begin{lemma}\label{lem:d3}
$\rank\partial_3=n$ and every invariant factor of $\partial_3$ is $1$.
\end{lemma}

\begin{proof}
The $n$ generators of $C_3(D)$ are the squares $Q_j$. For $n\ge5$ any two distinct squares meet in at
most two vertices ($Q_j\cap Q_{j+1}$ is a rung, and $Q_j\cap Q_{j'}=\emptyset$ otherwise), so no
triple lies in two squares and the $n$ columns of $\partial_3$ have pairwise disjoint supports with
entries $\pm1$. Choosing one triple from each square gives an $n\times n$ minor equal to $\pm1$;
the rank is therefore $n$.
\end{proof}

\begin{lemma}\label{lem:d1}
$\rank\partial_1=2n-1$ and every invariant factor of $\partial_1$ is $1$.
\end{lemma}

\begin{proof}
$C_1(D)$ has all $\binom{2n}{2}$ pairs as generators, so $\partial_1$ is the incidence matrix of the
complete graph $K_{2n}$. Its image lies in the sum-zero sublattice of $\Zm^{2n}$, so
$\rank\partial_1\le 2n-1$; the $2n-1$ columns indexed by a spanning star give a $(2n-1)\times(2n-1)$
minor $\pm1$ after deleting the row of the star's centre.
\end{proof}

\begin{lemma}\label{lem:d2}
$\rank\partial_2=5n$ and every invariant factor of $\partial_2$ is $1$.
\end{lemma}

\begin{proof}
\emph{Upper bound.} $\operatorname{im}\partial_3\subseteq\ker\partial_2$ and
$\rank\partial_3=n$ by Lemma~\ref{lem:d3}, so $\dim\ker\partial_2\ge n$ and
$\rank\partial_2\le 6n-n=5n$.

\emph{A $5n\times5n$ minor equal to $\pm1$.} Select the $5n$ columns
\[
 \mathcal C=\{T_1(v):v\in V\}\cup\{T_3(v):v\in V\}\cup\{T_2((1,j)):j\in\Zm_n\},
 \quad |\mathcal C|=5n,
\]
and the $5n$ rows
\[
 p_A(r,j)=\{(r,j{-}1),(r,j{+}1)\},\qquad
 p_B(r,j)=\{(r,j{+}1),(1{-}r,j)\},
\]
\[
 p_C(j)=\{(1,j{-}1),(1,j)\},
\]
for $r\in\{0,1\}$, $j\in\Zm_n$: respectively $2n$ same-rail pairs at distance $2$, $2n$ diagonal
pairs, and the $n$ rail edges of row~$1$. For $n\ge5$ these $5n$ pairs are distinct and of three
distinct types, so they are $5n$ distinct rows.

We repeatedly use one observation. If two vertices of a triple $T\in D$ are non-adjacent then, $T$
being an induced $P_3$, they are its two endpoints and its third vertex is a common neighbour of them.
So a non-adjacent pair $p$ lies in exactly as many triples of $D$ as the two vertices of $p$ have
common neighbours. For $n\ge5$: the pair $p_A(r,j)$ has the single common neighbour $(r,j)$, so it
lies only in $T_1((r,j))$; the pair $p_B(r,j)$ has the two common neighbours $(r,j)$ and $(1-r,j+1)$,
so it lies only in $T_3((r,j))$ and $T_2((1{-}r,j{+}1))$. (At $n=4$ the first statement fails, since
$(r,j-1)$ and $(r,j+1)$ then have two common neighbours; this is the second place the argument uses
$n\ge5$.) For the adjacent pair $p_C(j)$, the triples of $D$ containing it are those obtained by
adjoining a neighbour of one endpoint that is not adjacent to the other, namely
$T_1((1,j{-}1))$, $T_1((1,j))$, $T_3((1,j{-}1))$ and $T_2((1,j))$.

Now delete rows and columns from the $5n\times5n$ submatrix in four stages, at each stage using a row
whose only surviving nonzero entry is $\pm1$, so that the determinant is unchanged up to sign:
\begin{enumerate}
\item the $2n$ rows $p_A(r,j)$: each meets $\mathcal C$ only in $T_1((r,j))$. Delete them with the
 $2n$ columns $T_1(v)$.
\item the $n$ rows $p_B(1,j)$: each lies in exactly two triples of $D$, namely $T_3((1,j))$ and
 $T_2((0,j{+}1))$, and the latter is not in $\mathcal C$. Delete them with the $n$ columns
 $T_3((1,j))$.
\item the $n$ rows $p_C(j)$: of the four triples containing $p_C(j)$, three were deleted in stages 1
 and 2, leaving $T_2((1,j))$. Delete them with the $n$ columns $T_2((1,j))$.
\item the $n$ rows $p_B(0,j)$: each meets $\mathcal C$ in $T_3((0,j))$ and $T_2((1,j{+}1))$, and
 the latter was deleted in stage~3. Delete them with the $n$ columns $T_3((0,j))$.
\end{enumerate}
Nothing is left, so the minor is $\pm1$ and $\rank\partial_2\ge5n$.
\end{proof}

\subsection{The homology of \texorpdfstring{$D$}{D}, and of \texorpdfstring{$\Dl_3$}{Delta 3}}

\begin{proposition}\label{prop:dualhom}
For $n\ge5$, $\widetilde H_i(D;\Zm)=0$ for $i\ne1$, and $\widetilde H_1(D;\Zm)\cong\Zm^{2n^2-8n+1}$.
\end{proposition}

\begin{proof}
By Lemmas \ref{lem:d3}--\ref{lem:d2} every invariant factor of $\partial_1$, $\partial_2$ and
$\partial_3$ equals~$1$, and $\partial_0$ is the augmentation, which is surjective with its single
invariant factor~$1$. Since $\operatorname{im}\partial_{i+1}\subseteq\ker\partial_i$ and
$\ker\partial_i$ is free, the torsion subgroup of $\widetilde H_i(D;\Zm)$ is the direct sum of the
$\Zm/d_j$ over the invariant factors $d_j$ of $\partial_{i+1}$, hence trivial. So every
$\widetilde H_i(D;\Zm)$ is free and it suffices to compute ranks:
\[
\begin{aligned}
 \widetilde H_0&: && (2n-1)-\rank\partial_1 &&= (2n-1)-(2n-1)=0,\\
 \widetilde H_1&: && \tbinom{2n}{2}-\rank\partial_1-\rank\partial_2
   &&= n(2n-1)-(2n-1)-5n=2n^2-8n+1,\\
 \widetilde H_2&: && 6n-\rank\partial_2-\rank\partial_3 &&= 6n-5n-n=0,\\
 \widetilde H_3&: && n-\rank\partial_3 &&= n-n=0 .
\end{aligned}
\]
(For $\widetilde H_0$ note $\ker\partial_0$ has rank $2n-1$.) All other $\widetilde H_i$ vanish
because $\dim D=3$.
\end{proof}

\begin{proof}[Proof of the homology statement of Theorem~\ref{thm:main}]
By \eqref{eq:AD} with $N=2n$ and $K=\Dl_3(\CL_n)$, $K^{*}=D$, Proposition~\ref{prop:dualhom} gives
$\widetilde H^{\,2n-i-3}(\Dl_3;\Zm)\cong\widetilde H_i(D;\Zm)$, so
$\widetilde H^{\,j}(\Dl_3;\Zm)=0$ for every $j\ne2n-4$ and
$\widetilde H^{\,2n-4}(\Dl_3;\Zm)\cong\Zm^{2n^2-8n+1}$.

Now transport to homology. $\Dl_3$ is a finite complex, so all $\widetilde H_i$ are finitely
generated, and the universal coefficient theorem gives
$\widetilde H^{\,j}\cong\operatorname{Hom}(\widetilde H_j,\Zm)\oplus\operatorname{Ext}(\widetilde
H_{j-1},\Zm)$, that is, the free part of $\widetilde H_j$ direct sum the torsion of
$\widetilde H_{j-1}$. From $\widetilde H^{\,j}=0$ for $j\ne2n-4$ we get: the free part of
$\widetilde H_j$ is $0$ for every $j\ne2n-4$, and the torsion of $\widetilde H_{j-1}$ is $0$ for
every $j\ne2n-4$, i.e.\ $\widetilde H_i$ is torsion-free for every $i\ne2n-5$. At $j=2n-4$ the group
$\widetilde H^{\,2n-4}\cong\Zm^{2n^2-8n+1}$ is free, and it contains the torsion of
$\widetilde H_{2n-5}$ as a direct summand, so that torsion is trivial too. Combining, every
$\widetilde H_i$ with $i\ne2n-4$ is torsion-free with zero free part, hence zero, and
$\widetilde H_{2n-4}\cong\Zm^{2n^2-8n+1}$.
\end{proof}

\subsection{The homotopy type}

\begin{lemma}\label{lem:skel}
For $n\ge5$ every subset of $V$ of size at most $3$ is a face of $\Dl_3(\CL_n)$; that is,
$\Dl_3(\CL_n)$ contains the full $2$-skeleton of the simplex on $V$.
\end{lemma}

\begin{proof}
Let $|\sigma|\le3$, so $|V\setminus\sigma|\ge2n-3\ge7$. Pick $u\in V\setminus\sigma$. Among the
$\ge2n-4\ge6$ other vertices of $V\setminus\sigma$, at most $3$ are neighbours of $u$, so at least
$2n-7\ge3$ are not. Take two of them, $w_1,w_2$: the triple $\{u,w_1,w_2\}$ spans at most one edge
and so induces a disconnected subgraph. By \eqref{eq:facechar}, $\sigma\in\Dl_3(\CL_n)$.
\end{proof}

\begin{proof}[Proof of the homotopy statement of Theorem~\ref{thm:main}]
Let $X=\Dl_3(\CL_n)$ with $n\ge5$, let $d=2n-4\ge6$ and $r=2n^2-8n+1$. By Lemma~\ref{lem:skel} the
$2$-skeleton of $X$ is the $2$-skeleton of a simplex, which is simply connected; since $\pi_1$
depends only on the $2$-skeleton, $X$ is simply connected. By the homology statement
$\widetilde H_*(X;\Zm)$ is free and concentrated in the single degree $d$.

We first record that $X$ is $(d-1)$-connected, which is what the Hurewicz theorem needs. Suppose
$X$ is $(j-1)$-connected for some $j$ with $2\le j\le d-1$; since $j\ge2$, Hurewicz gives
$\pi_j(X)\cong\widetilde H_j(X;\Zm)=0$, and hence $X$ is $j$-connected. As $X$ is simply connected,
that is, $1$-connected, induction from $j=2$ gives $(d-1)$-connectedness. Hurewicz applied once more,
now in degree $d$, makes $\pi_d(X)\to H_d(X)\cong\Zm^{r}$ an isomorphism. Choose maps $S^{d}\to X$
representing a basis and let $f:\bigvee_{i=1}^{r}S^{d}\to X$ be the resulting map. Then $f$ induces an
isomorphism on integral homology in every degree, and both spaces are simply connected CW complexes,
so $f$ is a homotopy equivalence by the homology form of Whitehead's theorem \cite{Hatcher}. Note
$r=2n^2-8n+1\ge11>0$ for $n\ge5$, so the wedge is non-empty.
\end{proof}

\section{Scope, and verification}\label{sec:verify}

Part~(i) of Conjecture~\ref{conj:csv}, which concerns $n=4$ and $m\ge2$, remains open; nothing here
bears on $m\ge3$. The hypothesis $n\ge5$ is used twice, in Lemma~\ref{lem:fvec} and in
Lemma~\ref{lem:d2}, and part~(ii) is false at $n=4$: an exact integral computation gives
$\widetilde H_3(\Dl_3(\CL_4);\Zm)\cong\Zm$ and $\widetilde H_4(\Dl_3(\CL_4);\Zm)\cong\Zm^{4}$, and
zero elsewhere.

The accompanying program \texttt{verify.py} (Python 3.9+, standard library only, exact integer
arithmetic, no randomness) re-derives from the definitions: the graph invariants of $\CL_n$ and the
$f$-vector of $D$ for $n=4,\dots,40$, including the absence of a $5$-set and the value $f_3=6$ at
$n=4$ (Lemma~\ref{lem:fvec}); the exact integral homology of $\Dl_3(\CL_n)$ for $n=5,\dots,24$,
together with the boundary ranks $n$, $5n$, $2n-1$ of Lemmas \ref{lem:d3}--\ref{lem:d2} and the
vanishing of all invariant factors above $1$; the four-stage deletion of Lemma~\ref{lem:d2},
executed as a stripping certificate on the $5n\times5n$ submatrix itself for $n=5,\dots,24$; and the
same homology computed \emph{directly} from the full complex $\Dl_3$, with no duality and no relative
complex, at $n=4,5,6,7$. Its recorded output is \texttt{verify.output.txt} ($45$ checks, all
passing). Among the things the program states it does not cover are $n\ge25$, any $m\ge3$, the
homotopy-theoretic step, and any literature search; for those, and for all $n\ge5$, the paper rests
on the proof above.

\begin{thebibliography}{9}

\bibitem{BayerI}
M.~Bayer, M.~Denker, M.~Jeli\'c Milutinovi\'c, R.~Rowlands, S.~Sundaram, and L.~Xue,
\emph{Topology of cut complexes of graphs},
SIAM J. Discrete Math. \textbf{38} (2024), no.~2, 1630--1675;
\href{https://arxiv.org/abs/2304.13675}{arXiv:2304.13675}.

\bibitem{BayerII}
M.~Bayer, M.~Denker, M.~Jeli\'c Milutinovi\'c, R.~Rowlands, S.~Sundaram, and L.~Xue,
\emph{Topology of cut complexes II},
SIAM J. Discrete Math. \textbf{39} (2025), no.~2, 1123--1157.
\href{https://doi.org/10.1137/24M1676077}{doi:10.1137/24M1676077};
\href{https://arxiv.org/abs/2407.08158}{arXiv:2407.08158}.
The result used here is the one on graphs with no cycle of length $\le k$ and at least one cycle of
length $k+1$; it is stated in full in Section~\ref{sec:result} above.

\bibitem{BjornerTancer}
A.~Bj\"orner and M.~Tancer,
\emph{Combinatorial Alexander duality---a short and elementary proof},
Discrete Comput. Geom. \textbf{42} (2009), no.~4, 586--593.
\href{https://doi.org/10.1007/s00454-008-9102-x}{doi:10.1007/s00454-008-9102-x}.

\bibitem{CSV}
P.~Chauhan, S.~Shukla, and K.~Vinayak,
\emph{Total 2-cut complexes of powers of cycle graphs and Cartesian products of certain graphs},
\href{https://arxiv.org/abs/2512.04486}{arXiv:2512.04486v1}, 2025.
The conjecture settled here is the last numbered conjecture of that paper.

\bibitem{Hatcher}
A.~Hatcher,
\emph{Algebraic Topology},
Cambridge University Press, 2002.

\end{thebibliography}

\end{document}
